\setcounter{secnumdepth}{-1} % Esto evita que la introducción lleve número de capítulo (Ej: "Capítulo 0")

\chapter*{Introducción}
\addcontentsline{toc}{chapter}{Introducción}

La industria vitivinícola es uno de los sectores agroalimentarios de mayor impacto 
económico y cultural a nivel mundial. En España, con más de 2,6 millones de hectáreas 
de viñedo, representa uno de los pilares de la economía agraria y una seña de identidad 
nacional reconocida internacionalmente. La calidad del vino, entendida como la conformidad 
del producto con los estándares de su Denominación de Origen y las expectativas del 
consumidor, es el principal activo competitivo de cualquier bodega.

Sin embargo, garantizar esa calidad sigue siendo un desafío para el sector. Los métodos 
tradicionales de certificación, basados en análisis de laboratorio y en la evaluación 
sensorial de enólogos, son eficaces pero costosos, lentos y difíciles de escalar cuando 
el volumen de producción es elevado. Una bodega mediana puede necesitar evaluar miles de 
muestras por campaña, lo que supone una carga operativa y económica significativa.

Ante este escenario, la Inteligencia Artificial ofrece una oportunidad real para transformar 
la forma en que las bodegas gestionan la calidad de su producción. Este Hito 1 propone 
el diseño de un sistema inteligente capaz de clasificar la variedad de un vino a partir 
de sus propiedades fisicoquímicas, con el objetivo de apoyar al equipo técnico de la bodega 
en la toma de decisiones sobre la calidad de forma objetiva y reproducible.